\subsection{boundary-only faithful carrier：粘连缺陷、prime-log 留数与无限秩宿主}

正文已经分别给出：subfull 维 regime 中边界 Gödel 维的相图、精确边界码位长与条件 Kolmogorov 下界对粘连缺陷的读取、prime-log 校准的精确 coboundary 化、其对 bulk 谱压的不变性、以及 RH 候选算子在 prime-register 侧的可数无限秩障碍。
把这些结果再压缩为同一条边界承载原则，可得到如下两条派生后果。

\begin{theorem}[subfull 维 regime 中边界复杂度完全等于粘连缺陷]\label{thm:derived-subfull-boundary-complexity-equals-adhesion}
\leanverified{paper\_derived\_subfull\_boundary\_complexity\_equals\_adhesion}
设 \(A\subset I^n\) 为紧集，且
\[
d_G(A)<n.
\]
则边界复杂度的三种读法完全同一：
\begin{enumerate}
\normalfont
  \item 几何指数满足
  \[
  d_{\partial G}(A)=\gamma_{\mathrm{adh}}(A);
  \]
  \item 精确边界 Gödel 码位长
  \[
  L_{\partial,m}(A):=\left\lceil\log_2 H_m(U_m(A))\right\rceil
  \]
  满足
  \[
  \limsup_{m\to\infty}\frac{\log L_{\partial,m}(A)}{m\log 2}
  =
  \gamma_{\mathrm{adh}}(A);
  \]
  \item 若沿无穷子序列 \(m_j\to\infty\) 有
  \[
  K\!\left(H_{m_j}(U_{m_j}(A))\,\middle|\,m_j\right)\ge 2^{\kappa m_j},
  \]
  则
  \[
  \gamma_{\mathrm{adh}}(A)\ge \kappa.
  \]
\end{enumerate}
若再有紧集 \(A,B\subset I^n\) 满足
\[
\operatorname{dist}(A,B)>0,
\qquad
\max\{d_G(A),d_G(B)\}<n,
\]
则
\[
\gamma_{\mathrm{adh}}(A\cup B)
=
\max\{\gamma_{\mathrm{adh}}(A),\gamma_{\mathrm{adh}}(B)\}.
\]
更进一步，若体相复杂度强到足以迫使边界 Gödel 维饱和，则在 \(d_G(A)<n\) 的体制下必有
\[
\gamma_{\mathrm{adh}}(A)=n.
\]
因此，在非满维 regime 中，边界复杂度不再由体积厚度项主导，而被纯粘连亏损完整接管。
\end{theorem}

\begin{proof}
当 \(d_G(A)<n\) 时，定理
\ref{thm:conclusion-boundary-godel-adhesion-defect-phase-diagram}
直接给出
\[
d_{\partial G}(A)=\gamma_{\mathrm{adh}}(A).
\]
精确边界码位长指数的等式则是定理
\ref{thm:conclusion-boundary-godel-bitlength-exponent-equals-adhesion}
的内容。条件 Kolmogorov 下界推出 \(\gamma_{\mathrm{adh}}(A)\ge \kappa\) 则由定理
\ref{thm:conclusion-boundary-godel-kolmogorov-lowerbound-forces-adhesion}
给出。

对正距并集，极大律正是定理
\ref{thm:conclusion-boundary-godel-separated-union-adhesion-maxlaw}
的内容。最后，若体相算法复杂度达到使边界 Gödel 维饱和的 regime，则定理
\ref{thm:conclusion-boundary-kolmogorov-saturation-forces-godel-saturation}
与推论
\ref{cor:conclusion-boundary-kolmogorov-saturation-implies-maximal-adhesion}
给出
\[
d_{\partial G}(A)=n,
\qquad
\gamma_{\mathrm{adh}}(A)=n
\]
在 \(d_G(A)<n\) 的设定下同时成立。证毕。
\end{proof}

\begin{theorem}[prime-log 校准不改 bulk 谱压，而 faithful 宿主必为可数无限秩]\label{thm:derived-prime-log-boundary-only-current-countable-host}
\leanverified{paper\_derived\_prime\_log\_boundary\_only\_current\_countable\_host}
固定 \(B=2^L\) 与 \(\alpha>0\)，记
\[
c:=2\alpha\log B,
\qquad
\psi_\alpha(j):=y_jw_j-c.
\]
则以下各条同时成立：
\begin{enumerate}
\normalfont
  \item centered prime-log 势是精确 coboundary：
  \[
  \psi_\alpha(j)=\delta_{j-1}-\delta_j;
  \]
  \item 对任意 \(q\ge 1\) 与 \(n\ge 1\)，prime-log 扭曲后的 \(n\)-步射影算子满足
  \[
  \mathcal L_q^{(\alpha,n)}=e^{-\delta_n}\mathcal L_q^n,
  \qquad
  \lim_{n\to\infty}\frac1n\log\|\mathcal L_q^{(\alpha,n)}\|
  =
  \log\rho(\mathcal L_q),
  \]
  因而 \(\alpha\) 只产生 boundary twist，而不改变 bulk 谱压；
  \item Abel--Mellin 边界流满足
  \[
  \mathcal D_\alpha(s)=c\,\zeta(s)+\mathcal J_\alpha(s)
  \qquad (\Re s>1),
  \]
  其中 \(\mathcal J_\alpha\) 在 \(\Re s>0\) 上全纯，故 \(\alpha\) 的全部非全纯依赖严格压缩到
  \[
  \CC\cdot \zeta(s)
  \]
  这一维边界方向上；
  \item 若
  \[
  M_\alpha=
  \left\{
  \prod_{p\in\mathcal P_\alpha}p^{a_p}:\ 
  a_p\in\ZZ_{\ge 0},\ a_p=0\ \text{仅除有限多个 }p\text{ 外}
  \right\},
  \]
  则任何 faithful 地记录 \(M_\alpha\) 并恢复边界流 \(\Psi_\alpha\) 的候选 RH 宿主，都满足
  \[
  \ell_{\mathrm{loc}}(M_\alpha)=\aleph_0.
  \]
\end{enumerate}
因此，prime-log 校准的 \(\alpha\)-依赖只能作为一维边界留数进入，而与之兼容的 faithful RH 宿主不可能闭合在任何有限 local-axis、有限秩加法账本或其 Pontryagin 对偶的有限维 Lie 模型之内。
\end{theorem}

\begin{proof}
第一条正是定理
\ref{thm:conclusion-prime-log-exact-coboundary}
的内容。将其代入 \(n\)-步路径乘积，便得到第二条中的
\[
\mathcal L_q^{(\alpha,n)}=e^{-\delta_n}\mathcal L_q^n,
\]
而 bulk 谱压不变性则由定理
\ref{thm:conclusion-prime-log-boundary-twist-pressure-invariance}
直接给出。

第三条是定理
\ref{thm:conclusion-prime-log-abel-mellin-pole-rigidity}
与推论
\ref{cor:conclusion-prime-log-alpha-singular-rank-one}
的合并表述：\(\mathcal J_\alpha\) 在 \(\Re s>0\) 上全纯，而所有 singular \(\alpha\)-dependence 都落在 \(c\,\zeta(s)\) 这一维方向上。

最后，faithful 宿主的可数无限秩障碍正是定理
\ref{thm:conclusion-prime-log-rh-operator-countable-prime-rank}
的内容。于是 \(\alpha\) 不能作为 bulk Hamiltonian deformation 进入，只能作为边界留数/边界质量进入，而 faithful 宿主也必然是可数无限秩。证毕。
\end{proof}

\begin{remark}[boundary-only faithful carrier 原理]\label{rem:derived-boundary-only-faithful-carrier}
定理 \ref{thm:derived-subfull-boundary-complexity-equals-adhesion}
与定理 \ref{thm:derived-prime-log-boundary-only-current-countable-host}
共同表明：在本文所讨论的两类看似无关的界面上，真正不可压缩的量都不是 bulk 指数，而是某种最小 faithful 的边界载荷。其一，在 subfull 维 regime 中，这一载荷是纯粘连缺陷 \(\gamma_{\mathrm{adh}}\)；其二，在 prime-log 校准链中，这一载荷是一维边界留数 \(c\,\zeta(s)\) 及其所强迫的可数 prime-register 宿主。故“boundary-only carrier”并非偶然修辞，而是同一最小 faithful 原理在几何复杂度与谱实现两端的平行显现。
\end{remark}

\endinput
